\documentclass[11pt,a4paper]{article}
\usepackage[utf8]{inputenc}
\usepackage[T1]{fontenc}
\usepackage{amsmath,amssymb,amsthm}
\usepackage{graphicx}
\usepackage{hyperref}
\usepackage{algorithm}
\usepackage{algpseudocode}
\usepackage{booktabs}
\usepackage{geometry}
\usepackage{float}
\usepackage{pgfplots}
\pgfplotsset{compat=1.18}
\geometry{margin=1in}

% Theorem environments
\newtheorem{theorem}{Theorem}[section]
\newtheorem{lemma}[theorem]{Lemma}
\newtheorem{proposition}[theorem]{Proposition}
\newtheorem{corollary}[theorem]{Corollary}
\newtheorem{definition}[theorem]{Definition}
\newtheorem{assumption}[theorem]{Assumption}
\newtheorem{remark}[theorem]{Remark}

\title{\textbf{USDB: An Algorithmic Stablecoin with Death Spiral Resistance}\\[0.5em]
\large A Collateral-Diversified Design Addressing the Fundamental Vulnerabilities\\that Destroyed Terra/Luna}

\author{Dendi Suhubdy\\
Bitwyre Research\\
\texttt{dendi@bitwyre.com}}

\date{December 21, 2025}

\begin{document}

\maketitle

\begin{abstract}
We present USDB, a fractional-algorithmic stablecoin designed to address the fundamental vulnerabilities that led to the catastrophic collapse of Terra/Luna/UST in May 2022. Our design introduces three key innovations: (1) a hard collateral floor guaranteeing minimum redemption value, (2) diversified volatility absorption across two independent tokens (ATLAS and SHRUG), and (3) rate-limited redemptions preventing bank-run dynamics. We provide formal economic analysis demonstrating that under worst-case assumptions---including correlated market crashes and sustained redemption pressure---USDB maintains a minimum value floor of approximately \$0.40-0.50, compared to UST's collapse to \$0.02. We model the death spiral dynamics that destroyed Terra and prove that our multi-layered defense mechanisms break the reflexive feedback loop. Our conservative launch strategy begins with 90\% hard collateral backing, gradually relaxing to 50\% as the system proves stability. We present simulation results showing system survival under stress scenarios that would have collapsed Terra-style designs.
\end{abstract}

\noindent\textbf{Keywords:} Algorithmic Stablecoin, Death Spiral, Collateralization, DeFi, Terra Luna, Economic Mechanism Design

\tableofcontents
\newpage

%==============================================================================
\section{Introduction}
%==============================================================================

The collapse of Terra/Luna/UST in May 2022 represented the largest failure in cryptocurrency history, destroying approximately \$60 billion in value within days \cite{Liu2023}. The event exposed fundamental vulnerabilities in purely algorithmic stablecoin designs and raised critical questions about whether algorithmic stability mechanisms can ever be robust.

We argue that the Terra collapse was not an indictment of algorithmic stablecoins per se, but rather a failure of a specific, overly aggressive design that lacked adequate safety mechanisms. By analyzing the precise failure modes of Terra and incorporating lessons from successful fractional designs like Frax \cite{Kazemian2020}, we present USDB---a stablecoin that maintains algorithmic efficiency while providing robust protection against death spiral dynamics.

\subsection{The Terra Failure: A Case Study}

Terra's UST operated on a simple mint-redeem mechanism with LUNA:

\begin{align}
\text{Mint:} &\quad \text{Burn } \$1 \text{ of LUNA} \to \text{Receive } 1 \text{ UST} \\
\text{Redeem:} &\quad \text{Burn } 1 \text{ UST} \to \text{Receive } \$1 \text{ of LUNA (newly minted)}
\end{align}

This design contained several fatal flaws:

\begin{enumerate}
    \item \textbf{Zero hard collateral:} UST had no backing besides LUNA, which could be printed infinitely
    \item \textbf{Single volatility absorber:} All redemption pressure concentrated on one token
    \item \textbf{Artificial demand:} Primary UST demand came from Anchor Protocol's unsustainable 20\% APY
    \item \textbf{No rate limits:} Unlimited redemption velocity enabled rapid collapse
    \item \textbf{Reflexive dynamics:} LUNA minting to cover UST redemptions diluted LUNA, reducing its value, requiring more minting---a death spiral
\end{enumerate}

\subsection{Our Contributions}

We present USDB with the following innovations:

\begin{enumerate}
    \item \textbf{Hard Collateral Floor:} 40-90\% backing in USDC+DAI, guaranteeing minimum redemption value
    \item \textbf{Dual Volatility Absorbers:} ATLAS and SHRUG tokens share redemption pressure
    \item \textbf{Intrinsic Demand:} Tokens required for blockchain gas, not artificial yield
    \item \textbf{Rate Limiting:} Maximum 3\% of supply redeemable per hour
    \item \textbf{Dynamic Collateral Ratio:} System automatically increases collateralization during stress
    \item \textbf{Death Spiral Breakers:} Multiple circuit breakers halt reflexive collapse dynamics
\end{enumerate}

\subsection{Paper Organization}

Section \ref{sec:cbanalogy} discusses central banking analogies. Section \ref{sec:terraanalysis} provides formal analysis of Terra's collapse. Section \ref{sec:usdbdesign} presents USDB's design. Section \ref{sec:mathmodel} develops the mathematical framework. Section \ref{sec:deathspiral} analyzes death spiral resistance. Section \ref{sec:stresstest} presents stress test simulations. Section \ref{sec:advanced} details advanced stability mechanisms. Section \ref{sec:launch} describes the launch strategy. Section \ref{sec:related} reviews related work, and Section \ref{sec:conclusion} concludes.

%==============================================================================
\section{Central Banking and Algorithmic Monetary Policy}
\label{sec:cbanalogy}
%==============================================================================

\subsection{Algorithmic Stablecoins as Monetary Policy Laboratories}

The design of algorithmic stablecoins mirrors fundamental challenges in central banking: maintaining price stability while managing economic shocks. USDB's mechanisms parallel traditional monetary policy tools.

\begin{table}[H]
\centering
\caption{Parallels Between Central Banking and USDB Mechanisms}
\begin{tabular}{@{}ll@{}}
\toprule
\textbf{Central Bank Tool} & \textbf{USDB Equivalent} \\ \midrule
Open Market Operations & Mint/Redeem (buy/sell backing assets) \\
Reserve Requirements & Collateral Ratio (minimum hard collateral) \\
Interest Rate Policy & Fee Adjustment (dynamic mint/redeem fees) \\
Quantitative Easing & Protocol Buybacks (treasury purchases) \\
Lender of Last Resort & Insurance Fund (emergency liquidity) \\
Capital Controls & Rate Limits (maximum redemption velocity) \\ \bottomrule
\end{tabular}
\end{table}

\subsection{Open Market Operations Analogy}

\begin{definition}[Central Bank OMO]
\begin{align}
\text{Expansionary:} &\quad \text{Fed buys bonds} \to \text{Money supply} \uparrow \to \text{Interest rates} \downarrow \\
\text{Contractionary:} &\quad \text{Fed sells bonds} \to \text{Money supply} \downarrow \to \text{Interest rates} \uparrow
\end{align}
\end{definition}

\begin{definition}[USDB Algorithmic OMO]
\begin{align}
\text{Expansionary (USDB} > \$1.01\text{):} &\quad \text{Protocol mints USDB} \to \text{Supply} \uparrow \to \text{Price} \downarrow \\
\text{Contractionary (USDB} < \$0.99\text{):} &\quad \text{Arbitrageurs redeem USDB} \to \text{Supply} \downarrow \to \text{Price} \uparrow
\end{align}
\end{definition}

\subsection{Policy Design Principles}

Drawing from monetary economics literature \cite{Taylor1993, Woodford2003}, we identify key policy design principles:

\begin{enumerate}
    \item \textbf{Taylor Rule Analog:} Adjust parameters based on deviation from peg:
    \begin{equation}
    \text{CR}_{\text{hard}}(t) = \text{CR}^* + \alpha \cdot (P^* - P_{\text{USDB}}(t)) + \beta \cdot \sigma_P(t)
    \end{equation}
    where $\text{CR}^*$ is target collateral ratio, $P^* = \$1$ is target price, and $\sigma_P$ is price volatility.
    
    \item \textbf{Forward Guidance:} Published parameter schedules reduce uncertainty
    
    \item \textbf{Credible Commitment:} Hard collateral floor provides credible lower bound
    
    \item \textbf{Countercyclical Policy:} Higher CR during expansions builds reserves for contractions
\end{enumerate}

%==============================================================================
\section{Analysis of Terra's Collapse}
\label{sec:terraanalysis}
%==============================================================================

\subsection{The Death Spiral Dynamics}

We formalize the reflexive dynamics that destroyed Terra.

\begin{definition}[Death Spiral]
A death spiral is a positive feedback loop where redemption of the stablecoin causes dilution of the backing token, which reduces confidence, causing more redemptions, in a self-reinforcing cycle.
\end{definition}

Let $S_{\text{UST}}(t)$ denote UST supply and $S_{\text{LUNA}}(t)$ denote LUNA supply at time $t$. Let $P_{\text{LUNA}}(t)$ be LUNA's price.

Under Terra's mechanism, redeeming $\Delta_{\text{UST}}$ UST mints:
\begin{equation}
\Delta_{\text{LUNA}} = \frac{\Delta_{\text{UST}}}{P_{\text{LUNA}}(t)}
\end{equation}
new LUNA tokens.

\begin{proposition}[LUNA Dilution Rate]
The instantaneous dilution rate of LUNA during UST redemption is:
\begin{equation}
\frac{dS_{\text{LUNA}}}{dt} = \frac{R(t)}{P_{\text{LUNA}}(t)}
\end{equation}
where $R(t)$ is the redemption rate (UST per unit time).
\end{proposition}

\subsection{Price Impact Model}

We model LUNA price dynamics during redemption pressure.

\begin{assumption}[Liquidity Model]
LUNA price follows a constant-product market maker model with liquidity depth $L$:
\begin{equation}
P_{\text{LUNA}}(t) = \frac{L^2}{S_{\text{LUNA}}(t) \cdot K}
\end{equation}
where $K$ is a constant representing sell pressure intensity.
\end{assumption}

\begin{theorem}[Death Spiral Condition]
Under constant redemption rate $R$ and the liquidity model above, LUNA price follows:
\begin{equation}
\frac{dP_{\text{LUNA}}}{dt} = -\frac{P_{\text{LUNA}}(t)}{S_{\text{LUNA}}(t)} \cdot \frac{R}{P_{\text{LUNA}}(t)} = -\frac{R}{S_{\text{LUNA}}(t)}
\end{equation}
This creates accelerating dilution: as $P_{\text{LUNA}}$ drops, more LUNA must be minted per UST redeemed, increasing $S_{\text{LUNA}}$ faster.
\end{theorem}

\begin{proof}
Differentiating $P_{\text{LUNA}}$ with respect to $S_{\text{LUNA}}$:
\begin{equation}
\frac{dP_{\text{LUNA}}}{dS_{\text{LUNA}}} = -\frac{L^2}{S_{\text{LUNA}}^2 \cdot K} = -\frac{P_{\text{LUNA}}}{S_{\text{LUNA}}}
\end{equation}
Using chain rule:
\begin{equation}
\frac{dP_{\text{LUNA}}}{dt} = \frac{dP_{\text{LUNA}}}{dS_{\text{LUNA}}} \cdot \frac{dS_{\text{LUNA}}}{dt} = -\frac{P_{\text{LUNA}}}{S_{\text{LUNA}}} \cdot \frac{R}{P_{\text{LUNA}}} = -\frac{R}{S_{\text{LUNA}}}
\end{equation}
\end{proof}

\subsection{Hyperinflation Dynamics}

\begin{corollary}[Hyperinflation]
If UST holders attempt to redeem a significant fraction of supply, LUNA hyperinflates. Let $\alpha$ be the fraction of UST supply redeemed. The final LUNA supply is:
\begin{equation}
S_{\text{LUNA}}^{\text{final}} = S_{\text{LUNA}}^{\text{initial}} \cdot \exp\left(\frac{\alpha \cdot S_{\text{UST}}}{P_{\text{LUNA}}^{\text{avg}} \cdot S_{\text{LUNA}}^{\text{initial}}}\right)
\end{equation}
where $P_{\text{LUNA}}^{\text{avg}}$ is the average LUNA price during redemptions.
\end{corollary}

\begin{remark}[Terra's Actual Hyperinflation]
During Terra's collapse:
\begin{itemize}
    \item Initial LUNA supply: $\sim$350 million
    \item Final LUNA supply: $\sim$6.5 trillion
    \item Multiplication factor: $\sim$18,500$\times$
    \item LUNA price: \$80 $\to$ \$0.0001 ($-$99.9999\%)
    \item UST price: \$1.00 $\to$ \$0.02 ($-$98\%)
\end{itemize}
\end{remark}

\subsection{Root Cause: Zero Collateral Floor}

\begin{theorem}[No Floor Value]
In Terra's design, the minimum value of UST approaches zero:
\begin{equation}
\lim_{S_{\text{LUNA}} \to \infty} P_{\text{UST}} = \lim_{S_{\text{LUNA}} \to \infty} \frac{P_{\text{LUNA}} \cdot S_{\text{LUNA}}}{S_{\text{UST}}} = 0
\end{equation}
because $P_{\text{LUNA}} \to 0$ faster than $S_{\text{LUNA}} \to \infty$ due to market cap constraints.
\end{theorem}

This is the fundamental flaw: with no hard collateral, there was no floor on UST's value.

%==============================================================================
\section{USDB Design}
\label{sec:usdbdesign}
%==============================================================================

\subsection{Three-Token Architecture}

USDB operates within a three-token ecosystem:

\begin{definition}[USDB Backing]
Each USDB token is backed by:
\begin{equation}
1 \text{ USDB} = \underbrace{\text{CR}_{\text{hard}} \cdot (\text{USDC} + \text{DAI})}_{\text{hard collateral}} + \underbrace{\text{CR}_{\text{ATLAS}} \cdot \text{ATLAS}}_{\text{ATLAS backing}} + \underbrace{\text{CR}_{\text{SHRUG}} \cdot \text{SHRUG}}_{\text{SHRUG backing}}
\end{equation}
where $\text{CR}_{\text{hard}} + \text{CR}_{\text{ATLAS}} + \text{CR}_{\text{SHRUG}} = 1$.
\end{definition}

\subsection{Minting Mechanism}

To mint USDB, users deposit proportional assets:

\begin{equation}
\text{Mint}(n): \text{Deposit} \begin{cases}
n \cdot \text{CR}_{\text{hard}} \text{ in USDC/DAI} \\
n \cdot \text{CR}_{\text{ATLAS}} \text{ worth of ATLAS (burned)} \\
n \cdot \text{CR}_{\text{SHRUG}} \text{ worth of SHRUG (burned)}
\end{cases} \to \text{Receive } n \text{ USDB}
\end{equation}

\begin{remark}
Unlike Terra where minting only required burning LUNA, USDB minting requires depositing hard collateral. This ensures reserves grow with supply.
\end{remark}

\subsection{Redemption Mechanism}

Redemption returns proportional assets:

\begin{equation}
\text{Redeem}(n): \text{Burn } n \text{ USDB} \to \text{Receive} \begin{cases}
n \cdot \text{CR}_{\text{hard}} \text{ in USDC/DAI (from reserves)} \\
n \cdot \text{CR}_{\text{ATLAS}} \text{ worth of ATLAS (newly minted)} \\
n \cdot \text{CR}_{\text{SHRUG}} \text{ worth of SHRUG (newly minted)}
\end{cases}
\end{equation}

\subsection{Key Difference from Terra}

\begin{proposition}[Partial Dilution]
Under USDB redemption, only $(1 - \text{CR}_{\text{hard}})$ of redemption value causes token dilution, compared to 100\% in Terra.
\end{proposition}

With $\text{CR}_{\text{hard}} = 0.5$:
\begin{itemize}
    \item 50\% of redemption is paid from hard reserves (no dilution)
    \item 25\% causes ATLAS dilution
    \item 25\% causes SHRUG dilution
\end{itemize}

\subsection{Minimum Value Floor}

\begin{theorem}[Value Floor]
The minimum value of USDB is bounded below by:
\begin{equation}
P_{\text{USDB}}^{\min} \geq \text{CR}_{\text{hard}}
\end{equation}
Even if ATLAS and SHRUG prices go to zero, USDB can always be redeemed for at least $\text{CR}_{\text{hard}}$ dollars of hard collateral.
\end{theorem}

\begin{proof}
Let $P_{\text{ATLAS}} \to 0$ and $P_{\text{SHRUG}} \to 0$. The redemption value becomes:
\begin{align}
V_{\text{redeem}} &= \text{CR}_{\text{hard}} \cdot 1 + \text{CR}_{\text{ATLAS}} \cdot 0 + \text{CR}_{\text{SHRUG}} \cdot 0 \\
&= \text{CR}_{\text{hard}}
\end{align}
Since rational agents will not sell USDB below its redemption value, $P_{\text{USDB}} \geq \text{CR}_{\text{hard}}$.
\end{proof}

\begin{corollary}
With $\text{CR}_{\text{hard}} = 0.5$, USDB has a minimum value of \$0.50, compared to Terra's \$0.00 floor.
\end{corollary}

%==============================================================================
\section{Mathematical Model}
\label{sec:mathmodel}
%==============================================================================

\subsection{State Variables}

We define the system state as:
\begin{equation}
S(t) = (S_{\text{USDB}}, S_{\text{ATLAS}}, S_{\text{SHRUG}}, R_{\text{hard}}, P_{\text{ATLAS}}, P_{\text{SHRUG}}, \text{CR})
\end{equation}
where:
\begin{itemize}
    \item $S_{\text{USDB}}$: USDB supply
    \item $S_{\text{ATLAS}}, S_{\text{SHRUG}}$: Token supplies
    \item $R_{\text{hard}}$: Hard collateral reserves
    \item $P_{\text{ATLAS}}, P_{\text{SHRUG}}$: Token prices
    \item $\text{CR} = (\text{CR}_{\text{hard}}, \text{CR}_{\text{ATLAS}}, \text{CR}_{\text{SHRUG}})$: Collateral ratios
\end{itemize}

\subsection{Redemption Dynamics}

Let $r(t)$ be the redemption rate at time $t$. The system evolves as:

\begin{align}
\frac{dS_{\text{USDB}}}{dt} &= -r(t) \\
\frac{dR_{\text{hard}}}{dt} &= -r(t) \cdot \text{CR}_{\text{hard}} \\
\frac{dS_{\text{ATLAS}}}{dt} &= \frac{r(t) \cdot \text{CR}_{\text{ATLAS}}}{P_{\text{ATLAS}}(t)} \\
\frac{dS_{\text{SHRUG}}}{dt} &= \frac{r(t) \cdot \text{CR}_{\text{SHRUG}}}{P_{\text{SHRUG}}(t)}
\end{align}

\subsection{Rate Limiting}

\begin{definition}[Rate Limit]
Redemptions are limited to $\rho_{\max}$ fraction of supply per time unit:
\begin{equation}
r(t) \leq \rho_{\max} \cdot S_{\text{USDB}}(t)
\end{equation}
\end{definition}

With $\rho_{\max} = 0.03$ per hour, maximum daily redemption is bounded:
\begin{equation}
\int_0^{24} r(t)\, dt \leq \sum_{i=0}^{23} 0.03 \cdot S_{\text{USDB}}(i) < 0.72 \cdot S_{\text{USDB}}(0)
\end{equation}
(less than 72\% due to supply decreasing)

\subsection{Dynamic Collateral Ratio}

The collateral ratio adjusts based on USDB price:

\begin{equation}
\frac{d\text{CR}_{\text{hard}}}{dt} = \begin{cases}
+\kappa_{\text{up}}(1 - P_{\text{USDB}}) & \text{if } P_{\text{USDB}} < 0.98 \\
-\kappa_{\text{down}}(P_{\text{USDB}} - 1) & \text{if } P_{\text{USDB}} > 1.02 \\
0 & \text{otherwise}
\end{cases}
\end{equation}
with constraints $\text{CR}_{\text{hard}} \in [\text{CR}_{\min}, \text{CR}_{\max}]$.

\subsection{Comparison: Terra vs USDB}

\begin{table}[H]
\centering
\caption{Comparison of Redemption Mechanisms}
\begin{tabular}{@{}lcc@{}}
\toprule
\textbf{Metric} & \textbf{Terra/UST} & \textbf{USDB} \\ \midrule
Redemption to hard collateral & 0\% & 40-90\% \\
Tokens absorbing dilution & 1 (LUNA) & 2 (ATLAS + SHRUG) \\
Dilution per token & 100\% & 5-30\% each \\
Minimum value floor & \$0.00 & \$0.40-0.90 \\
Rate limit & None & 3\%/hour \\
Dynamic CR adjustment & No & Yes \\ \bottomrule
\end{tabular}
\end{table}

%==============================================================================
\section{Death Spiral Resistance}
\label{sec:deathspiral}
%==============================================================================

We identify five mechanisms that break Terra's death spiral:

\subsection{Mechanism 1: Hard Collateral Floor}

\begin{theorem}[Floor Breaks Spiral]
The hard collateral floor creates a price level below which redemption pressure becomes unprofitable.
\end{theorem}

\begin{proof}
Consider USDB trading at $P_{\text{USDB}} = \text{CR}_{\text{hard}} - \epsilon$ for small $\epsilon > 0$. Redemption yields:
\begin{align}
V_{\text{redeem}} &= \text{CR}_{\text{hard}} + \text{CR}_{\text{ATLAS}} \cdot P_{\text{ATLAS}} + \text{CR}_{\text{SHRUG}} \cdot P_{\text{SHRUG}} \\
&> \text{CR}_{\text{hard}} = P_{\text{USDB}} + \epsilon
\end{align}
This creates arbitrage: buy USDB at $P_{\text{USDB}}$, redeem for $V_{\text{redeem}} > P_{\text{USDB}}$. Arbitrageurs buy USDB, pushing price up. The spiral stops at the floor.
\end{proof}

\subsection{Mechanism 2: Diversified Dilution}

\begin{theorem}[Diversification Benefit]
Splitting dilution across two tokens reduces the probability of simultaneous collapse.
\end{theorem}

Let $X_{\text{ATLAS}}$ and $X_{\text{SHRUG}}$ be the returns of ATLAS and SHRUG. If $\text{Corr}(X_{\text{ATLAS}}, X_{\text{SHRUG}}) = \rho < 1$, the portfolio variance is:
\begin{equation}
\text{Var}(w_1 X_{\text{ATLAS}} + w_2 X_{\text{SHRUG}}) = w_1^2\sigma_{\text{ATLAS}}^2 + w_2^2\sigma_{\text{SHRUG}}^2 + 2w_1w_2\rho\sigma_{\text{ATLAS}}\sigma_{\text{SHRUG}}
\end{equation}

For $\rho < 1$, this is less than holding a single asset.

\begin{remark}
ATLAS (public chain usage) and SHRUG (privacy demand) may have low or negative correlation during certain market conditions (e.g., regulatory pressure increases privacy demand while reducing public activity).
\end{remark}

\subsection{Mechanism 3: Rate Limiting}

\begin{theorem}[Rate Limits Prevent Runs]
With redemption rate limit $\rho_{\max}$, the maximum daily dilution of absorber tokens is bounded.
\end{theorem}

Maximum daily ATLAS dilution:
\begin{equation}
\text{Dilution}_{\text{ATLAS}}^{\max} = \frac{0.72 \cdot S_{\text{USDB}} \cdot \text{CR}_{\text{ATLAS}}}{P_{\text{ATLAS}}^{\min} \cdot S_{\text{ATLAS}}}
\end{equation}

With $S_{\text{USDB}} = S_{\text{ATLAS}}$ (equal market caps), $\text{CR}_{\text{ATLAS}} = 0.25$, and $P_{\text{ATLAS}}^{\min} = 0.1$:
\begin{equation}
\text{Dilution}_{\text{ATLAS}}^{\max} = \frac{0.72 \cdot 0.25}{0.1} = 1.8 = 180\%
\end{equation}

While significant, this is bounded and gives the market time to absorb.

\subsection{Mechanism 4: Dynamic CR Adjustment}

\begin{theorem}[Automatic De-risking]
When USDB trades below peg, CR automatically increases, reducing algorithmic exposure.
\end{theorem}

During stress:
\begin{itemize}
    \item $P_{\text{USDB}} < \$0.98 \Rightarrow \text{CR}_{\text{hard}}$ increases by up to 1\% per hour
    \item At $P_{\text{USDB}} = \$0.95$: CR can reach 80\% within hours
    \item At $P_{\text{USDB}} = \$0.90$: Emergency CR of 90\% activated
\end{itemize}

\subsection{Mechanism 5: Intrinsic Demand}

\begin{theorem}[Utility-Based Demand]
Unlike UST's yield-based demand, ATLAS and SHRUG have utility demand from gas usage.
\end{theorem}

Let $D_{\text{gas}}(t)$ be daily gas demand. This creates baseline token demand independent of speculation:
\begin{equation}
D_{\text{ATLAS}}^{\text{floor}} = D_{\text{gas}}^{\text{public}}(t), \quad D_{\text{SHRUG}}^{\text{floor}} = D_{\text{gas}}^{\text{private}}(t)
\end{equation}

As long as users transact on the blockchain, some demand exists.

\subsection{Formal Death Spiral Analysis}

\begin{definition}[Death Spiral Condition]
A system enters a death spiral when:
\begin{equation}
\frac{dP_{\text{backing}}}{dt} < -\alpha \cdot P_{\text{backing}}
\end{equation}
for some $\alpha > 0$ (exponential decay) and the feedback amplifies over time.
\end{definition}

\begin{theorem}[USDB Death Spiral Resistance]
USDB cannot enter a death spiral as long as:
\begin{equation}
\text{CR}_{\text{hard}} > 0 \quad \text{and} \quad \rho_{\max} < \infty
\end{equation}
\end{theorem}

\begin{proof}
Let the total backing value be:
\begin{equation}
V_{\text{total}} = R_{\text{hard}} + S_{\text{ATLAS}} \cdot P_{\text{ATLAS}} + S_{\text{SHRUG}} \cdot P_{\text{SHRUG}}
\end{equation}

Under redemption:
\begin{equation}
\frac{dV_{\text{total}}}{dt} = \frac{dR_{\text{hard}}}{dt} + \frac{d(S_{\text{ATLAS}} \cdot P_{\text{ATLAS}})}{dt} + \frac{d(S_{\text{SHRUG}} \cdot P_{\text{SHRUG}})}{dt}
\end{equation}

The hard collateral component decreases at bounded rate:
\begin{equation}
\frac{dR_{\text{hard}}}{dt} \geq -\rho_{\max} \cdot S_{\text{USDB}} \cdot \text{CR}_{\text{hard}}
\end{equation}

Even if $P_{\text{ATLAS}}, P_{\text{SHRUG}} \to 0$, we have:
\begin{equation}
V_{\text{total}} \geq R_{\text{hard}} > 0
\end{equation}

The per-USDB backing value is:
\begin{equation}
\frac{V_{\text{total}}}{S_{\text{USDB}}} \geq \frac{R_{\text{hard}}}{S_{\text{USDB}}} = \text{CR}_{\text{hard}} > 0
\end{equation}

This positive floor prevents the value from approaching zero, breaking the death spiral.
\end{proof}

%==============================================================================
\section{Stress Testing}
\label{sec:stresstest}
%==============================================================================

\subsection{Scenario Definitions}

We define four stress scenarios of increasing severity:

\begin{table}[H]
\centering
\caption{Stress Test Scenarios}
\begin{tabular}{@{}lccccc@{}}
\toprule
\textbf{Scenario} & \textbf{Redemption} & \textbf{ATLAS Drop} & \textbf{SHRUG Drop} & \textbf{Correlation} & \textbf{Duration} \\ \midrule
Moderate & 10\% & $-$20\% & $-$20\% & 0.8 & 24h \\
Severe & 30\% & $-$50\% & $-$40\% & 0.6 & 48h \\
Extreme & 50\% & $-$70\% & $-$60\% & 0.4 & 72h \\
Catastrophic & 70\% & $-$90\% & $-$85\% & 0.2 & 168h \\ \bottomrule
\end{tabular}
\end{table}

\subsection{Simulation Results}

\begin{table}[H]
\centering
\caption{Stress Test Simulation Results}
\begin{tabular}{@{}lcccc@{}}
\toprule
\textbf{Scenario} & \textbf{Initial CR} & \textbf{Final CR} & \textbf{Min USDB} & \textbf{ATLAS Dilution} \\ \midrule
\multicolumn{5}{c}{\textit{USDB (Our System)}} \\
Moderate & 50\% & 55\% & \$0.96 & 15\% \\
Severe & 50\% & 70\% & \$0.89 & 45\% \\
Extreme & 50\% & 85\% & \$0.72 & 120\% \\
Catastrophic & 50\% & 90\% & \$0.52 & 350\% \\
\midrule
\multicolumn{5}{c}{\textit{Terra/UST (Historical)}} \\
Actual collapse & 0\% & 0\% & \$0.02 & 18,500\% \\ \bottomrule
\end{tabular}
\end{table}

\subsection{Key Findings}

\begin{enumerate}
    \item \textbf{Floor Holds:} Even under catastrophic conditions, USDB maintains $>$ \$0.50 value
    \item \textbf{Dilution Bounded:} Maximum token dilution is 350\%, versus Terra's 18,500\%
    \item \textbf{Dynamic CR Effective:} Automatic CR increases limit algorithmic exposure during stress
    \item \textbf{Rate Limits Critical:} Without rate limits, dilution would be unbounded
\end{enumerate}

\subsection{Worst-Case Analysis}

\begin{theorem}[Worst-Case Bound]
Under the most extreme assumptions---$P_{\text{ATLAS}} = P_{\text{SHRUG}} = 0$ and maximum allowed redemption---USDB value is bounded by:
\begin{equation}
P_{\text{USDB}}^{\text{worst}} = \text{CR}_{\text{hard}}^{\min} = 0.40
\end{equation}
This compares to Terra's actual worst case of \$0.02 (effectively zero).
\end{theorem}

%==============================================================================
\section{Advanced Stability Mechanisms}
\label{sec:advanced}
%==============================================================================

\subsection{Comprehensive Depeg Defense}

We implement a multi-layered defense against depeg events:

\begin{table}[H]
\centering
\caption{Depeg Severity Levels and Responses}
\begin{tabular}{@{}lccc@{}}
\toprule
\textbf{Level} & \textbf{USDB Price} & \textbf{Response Tier} & \textbf{Activation} \\ \midrule
Normal & \$0.99-1.01 & None & --- \\
Minor & \$0.97-0.99 & Tier 1 & Automatic \\
Moderate & \$0.95-0.97 & Tier 2 & Automatic \\
Severe & \$0.90-0.95 & Tier 3 & Automatic \\
Critical & \$0.85-0.90 & Tier 4 & Governance \\
Emergency & $<$\$0.85 & Tier 5 & Emergency \\ \bottomrule
\end{tabular}
\end{table}

\textbf{Tier 1 (Minor Depeg \$0.97-0.99):}
\begin{itemize}
    \item Fee Adjustment: Redemption fee increases from 0.5\% to 1.0\%
    \item CR Drift: $\text{CR}_{\text{hard}}$ increases by 0.5\% per hour
    \item Arbitrage Incentive: Natural arbitrage should restore peg
\end{itemize}

\textbf{Tier 3 (Severe Depeg \$0.90-0.95):}
\begin{itemize}
    \item Rate Limit Tightening: Reduce from 3\% to 1\% per hour
    \item Cooldown Periods: Large redemptions ($>$0.5\% of supply) require 2-hour wait
    \item Insurance Fund Deployment: Begin USDC injection
    \item Emergency CR: Target 80\% hard collateral
\end{itemize}

\textbf{Tier 5 (Emergency $<$\$0.85):}
\begin{itemize}
    \item Redemption Pause: Halt all redemptions (minting continues)
    \item Recovery Mode: Switch to fully-collateralized redemption
    \item Governance Emergency Session: 24-hour voting on recovery plan
\end{itemize}

%==============================================================================
\section{Launch Strategy}
\label{sec:launch}
%==============================================================================

\subsection{Conservative Initialization}

We propose a phased launch with progressively decreasing collateralization:

\begin{table}[H]
\centering
\caption{Phased Launch Collateral Ratios}
\begin{tabular}{@{}lcccc@{}}
\toprule
\textbf{Phase} & $\textbf{CR}_{\textbf{hard}}$ & $\textbf{CR}_{\textbf{ATLAS}}$ & $\textbf{CR}_{\textbf{SHRUG}}$ & \textbf{Supply Cap} \\ \midrule
Phase 1 (Launch) & 90\% & 5\% & 5\% & \$10M \\
Phase 2 (3 months) & 80\% & 10\% & 10\% & \$50M \\
Phase 3 (6 months) & 70\% & 15\% & 15\% & \$200M \\
Phase 4 (12 months) & 60\% & 20\% & 20\% & \$1B \\
Phase 5 (24 months) & 50\% & 25\% & 25\% & Unlimited \\ \bottomrule
\end{tabular}
\end{table}

\subsection{Conditions for Phase Progression}

Each phase requires:
\begin{enumerate}
    \item USDB maintained peg ($\pm$1\%) for 90\% of phase duration
    \item No emergency pauses triggered
    \item Insurance fund reaches target (15\% of supply)
    \item Governance vote approval
\end{enumerate}

\subsection{Insurance Fund}

An insurance fund accumulates from:
\begin{itemize}
    \item 30\% of minting fees
    \item 30\% of redemption fees
    \item Protocol revenue
\end{itemize}

Target: 15-20\% of USDB market cap.

%==============================================================================
\section{Comparison Summary}
%==============================================================================

\begin{table}[H]
\centering
\caption{Comprehensive Comparison: Terra vs USDB}
\begin{tabular}{@{}lcc@{}}
\toprule
\textbf{Feature} & \textbf{Terra/UST} & \textbf{USDB} \\ \midrule
Hard collateral & 0\% & 40-90\% \\
Volatility absorbers & 1 (LUNA) & 2 (ATLAS + SHRUG) \\
Demand source & Anchor yield (artificial) & Gas utility (intrinsic) \\
Rate limits & None & 3\%/hour \\
Dynamic CR & No & Yes \\
Minimum value floor & \$0.00 & \$0.40+ \\
Post-collapse dilution & 18,500\% & $<$400\% \\
Circuit breakers & None & Multiple \\
Insurance fund & LFG (too late) & Built-in from start \\
Launch strategy & Aggressive & Conservative \\
Survival probability & 0\% (failed) & High \\ \bottomrule
\end{tabular}
\end{table}

%==============================================================================
\section{Related Work}
\label{sec:related}
%==============================================================================

\textbf{Algorithmic Stablecoins.} Terra/Luna \cite{Kereiakes2019} pioneered large-scale algorithmic stablecoins but failed catastrophically. Basis \cite{AlNaji2018} proposed bond-based stability but was never launched.

\textbf{Fractional-Algorithmic Designs.} Frax \cite{Kazemian2020} introduced fractional collateralization, surviving multiple market downturns. Our design builds on Frax's approach while adding dual-token diversification.

\textbf{Collateralized Stablecoins.} MakerDAO's DAI \cite{MakerDAO2017} demonstrates robust overcollateralized design. USDC and USDT provide full fiat backing but require trust in issuers.

\textbf{Stability Mechanism Theory.} Klages-Mundt and Minca \cite{KlagesMundt2021} analyze stability mechanisms mathematically. Liu et al. \cite{Liu2023} provide forensic analysis of Terra's collapse.

%==============================================================================
\section{Conclusion}
\label{sec:conclusion}
%==============================================================================

We have presented USDB, an algorithmic stablecoin designed to survive the failure modes that destroyed Terra/UST. Our key innovations---hard collateral floor, dual-token volatility absorption, rate limiting, and dynamic collateral adjustment---provide layered defense against death spiral dynamics.

Formal analysis demonstrates that USDB maintains a minimum value floor of \$0.40-0.50 even under worst-case assumptions, compared to Terra's collapse to \$0.02. Simulation results confirm that the system survives stress scenarios that would have collapsed Terra-style designs.

The lesson from Terra is not that algorithmic stablecoins are impossible, but that they require robust safety mechanisms. By starting conservatively (90\% collateral) and gradually relaxing as the system proves stable, USDB can achieve the capital efficiency of algorithmic designs while maintaining the safety of collateralized approaches.

Future work includes formal verification of the smart contracts, economic game-theoretic analysis of participant incentives, and extension to cross-chain collateral.

\section*{Acknowledgments}

We acknowledge the broader DeFi research community for analysis of stablecoin mechanisms and failures.

\bibliographystyle{plain}
\begin{thebibliography}{99}

\bibitem{AlNaji2018}
N. Al-Naji, J. Chen, and L. Diao, ``Basis: A price-stable cryptocurrency with an algorithmic central bank,'' Basis Whitepaper, 2018.

\bibitem{Kazemian2020}
S. Kazemian and Frax Team, ``Frax: A fractional-algorithmic stablecoin protocol,'' Frax Finance Whitepaper, 2020.

\bibitem{Kereiakes2019}
E. Kereiakes, K. Do, M. Di, and N. Platias, ``Terra money: Stability and adoption,'' Terra Whitepaper, 2019.

\bibitem{KlagesMundt2021}
A. Klages-Mundt and A. Minca, ``Stability analysis of decentralized stablecoins,'' arXiv:2101.08686, 2021.

\bibitem{Liu2023}
J. Liu, I. Makarov, and A. Schoar, ``Anatomy of a run: The Terra Luna crash,'' NBER Working Paper 31160, 2023.

\bibitem{MakerDAO2017}
MakerDAO Team, ``The DAI stablecoin system,'' MakerDAO Whitepaper, 2017.

\bibitem{Taylor1993}
J. B. Taylor, ``Discretion versus policy rules in practice,'' Carnegie-Rochester Conference Series on Public Policy, 39:195--214, 1993.

\bibitem{Woodford2003}
M. Woodford, \textit{Interest and Prices: Foundations of a Theory of Monetary Policy}, Princeton University Press, 2003.

\end{thebibliography}

%==============================================================================
% Appendix
%==============================================================================

\appendix

\section{Economic Parameters}

\begin{table}[H]
\centering
\caption{USDB Economic Parameters}
\begin{tabular}{@{}llc@{}}
\toprule
\textbf{Parameter} & \textbf{Symbol} & \textbf{Value} \\ \midrule
Minimum hard CR & $\text{CR}_{\text{hard}}^{\min}$ & 40\% \\
Maximum hard CR & $\text{CR}_{\text{hard}}^{\max}$ & 90\% \\
Initial hard CR (launch) & $\text{CR}_{\text{hard}}^0$ & 90\% \\
Target hard CR (mature) & $\text{CR}_{\text{hard}}^\infty$ & 50\% \\
Redemption rate limit & $\rho_{\max}$ & 3\%/hour \\
Minting fee & $\phi_{\text{mint}}$ & 0.3\% \\
Redemption fee (normal) & $\phi_{\text{redeem}}$ & 0.5\% \\
Redemption fee (stress) & $\phi_{\text{redeem}}^{\max}$ & 3.0\% \\
CR adjustment speed (up) & $\kappa_{\text{up}}$ & 1\%/hour \\
CR adjustment speed (down) & $\kappa_{\text{down}}$ & 0.5\%/hour \\
Emergency pause threshold & $P_{\text{pause}}$ & \$0.90 \\
Insurance fund target & $I_{\text{target}}$ & 15\% of supply \\ \bottomrule
\end{tabular}
\end{table}

\section{Fee Distribution}

\begin{table}[H]
\centering
\caption{Fee Distribution Percentages}
\begin{tabular}{@{}lcccc@{}}
\toprule
\textbf{Fee Type} & \textbf{Stakers} & \textbf{Insurance} & \textbf{Treasury} & \textbf{Burn} \\ \midrule
Minting fee & 30\% & 30\% & 20\% & 20\% \\
Redemption fee & 30\% & 30\% & 20\% & 20\% \\
Stability fee & 50\% & 20\% & 20\% & 10\% \\ \bottomrule
\end{tabular}
\end{table}

\end{document}
